% §3 Method

\section{Method}
\label{sec:method}

We formalize unsupervised pre-inference routing (\S\ref{sec:method:formulation}), define \textbf{latent routing dimensions} and their operationalizations (\S\ref{sec:method:dimensions}), describe prefill-based signal extraction (\S\ref{sec:method:extraction}), and define evaluation metrics.
Experimental configuration: \S\ref{sec:setup}; empirical evidence: \S\ref{sec:characterization}; routing evaluation: \S\ref{sec:routing}.

\subsection{Routing formulation}
\label{sec:method:formulation}

Let $\mathcal{M}=\{M_1,\ldots,M_k\}$ be a fixed pool with known cost tiers \citep{hu2024routerbench,ding2024hybrid}.
For each query $q$, routing assigns one $M_i \in \mathcal{M}$ with objective \textbf{appropriate model selection}: route to the \textbf{weakest model expected to answer correctly}.

\textbf{Unsupervised pre-inference routing} uses query text and prefill probes before full generation, without routing supervision at extraction time \citep{guha2024smoothie,mu2025ragrouting}.
Prior work often ranks isolated statistics (e.g., entropy vs.\ margin).
We instead ask which \textbf{latent routing dimensions}---unobserved aspects of routing need---are measurable before generation and whether they predict \textbf{routing opportunity}.

\paragraph{Research hypotheses.}
\begin{align*}
  \text{RH1:}& \;\text{latent routing dimensions carry measurable predictive content for routing opportunity} \\
  \text{RH2:}& \;\text{dimensions encode different aspects of routing need} \\
  \text{RH3:}& \;\text{dimensions provide complementary predictive content beyond any single dimension} \\
  \text{RH4:}& \;\text{a calibrated routing policy can exploit available information.}
\end{align*}

\paragraph{Routing need (evaluation).}
For $(M_w,M_s)$, \textbf{routing opportunity} is $y(q,M_w)=0$ and $y(q,M_s)=1$ ($y^{\text{opp}}$).
Offline oracle labels $y(q,M_i)$ come from full greedy inference (\S\ref{sec:method:extraction}); they never enter signal extraction.
Four buckets---easy, opportunity, weak-only, too-hard---structure dimension-wise characterization.

\subsection{Latent routing dimensions}
\label{sec:method:dimensions}

A \textbf{latent routing dimension} is an aspect of routing need that is not observed directly---for example, how difficult the task is for the pool, how uncertain the weak model is, or whether escalation to the strong model is likely to help.
Each dimension is \textbf{operationalized} by a deterministic statistic computed from query text or prefill logits (Table~\ref{tab:dimensions}).

\paragraph{Scientific unit.}
The unit of analysis is the \textbf{dimension}, not the statistic.
Entropy, \texttt{piece\_count}, and $\Delta m_{\mathrm{gain}}$ are \emph{operationalizations}; the research question is whether the underlying dimension predicts opportunity.
Studies I--II therefore ask, for each dimension: \textbf{does this aspect of routing need predict $y^{\text{opp}}$?}
We report $\rho$ and AUROC as summaries of detectability---not as a feature leaderboard.

% Latent routing dimensions and their operationalizations (main analysis).
\begin{table}[t]
  \centering
  \small
  \setlength{\tabcolsep}{5pt}
  \begin{tabular}{@{}llp{0.42\linewidth}@{}}
    \toprule
    \textbf{Latent routing dimension} & \textbf{Operationalization} & \textbf{Measurement source} \\
    \midrule
    \textbf{Task difficulty} &
    $c(q)=$ \texttt{piece\_count} &
    Query text (model-independent) \\
    \addlinespace
    \textbf{Model uncertainty} &
    $H_w$ &
    Weak-model prefill logits \\
    \addlinespace
    \textbf{Model disagreement} &
    $\Delta H = H_w - H_s$ &
    Paired weak/strong logits \\
    \addlinespace
    \textbf{Escalation potential} &
    $\Delta m_{\mathrm{gain}} = m_s - m_w$ &
    Paired weak/strong logits \\
    \bottomrule
  \end{tabular}
  \caption{Latent routing dimensions and operationalizations studied in the main analysis. Dimensions are \emph{not} observed directly; each is measured through one CALIB-selected statistic. Appendix audits alternate operationalizations (e.g., $m_w$ for model uncertainty; Table~\ref{tab:md-variants}).}
  \label{tab:dimensions}
\end{table}


\paragraph{Operationalization details.}
\textbf{Task difficulty} is measured by $c(q)=$ \texttt{piece\_count}, selected on CALIB from query-text candidates (appendix Table~\ref{tab:d46}).
\textbf{Model uncertainty} uses weak-model entropy $H_w$ at the final prefill position \citep{he2023uncertainty,plaut2024chat}.
\textbf{Model disagreement} and \textbf{escalation potential} use paired weak/strong logits: $\Delta H = H_w - H_s$ and $\Delta m_{\mathrm{gain}} = m_s - m_w$.
Appendix Table~\ref{tab:md-variants} audits alternate operationalizations (e.g., $m_w$ as an alternative to $H_w$ for model uncertainty).

\subsection{Prefill-based signal extraction}
\label{sec:method:extraction}

Table~\ref{tab:pipeline} summarizes the analytical pipeline.
\textbf{Signal extraction} performs one prefill forward pass per pool member.
Richer measurements (trajectories, hidden states) would change operationalizations but not the dimensional abstraction.

% Analytical pipeline (orthogonal to dimension taxonomy in Table~\ref{tab:dimensions}).
\begin{table}[t]
  \centering
  \small
  \begin{tabular}{@{}clp{0.58\linewidth}@{}}
    \toprule
    & \textbf{Stage} & \textbf{Description} \\
    \midrule
    1 & Signal extraction &
    Query text; prefill logits at $t{=}T$ per pool member \\
    2 & Statistic computation (offline) &
    One representative per studied dimension (Table~\ref{tab:dimensions}) \\
    3 & Signal characterization &
    Routing-relevant information; cross-dimension structure; complementarity \\
    4 & Routing evaluation &
    Calibrated logistic policy; offline oracle; exploitation gap \\
    \bottomrule
  \end{tabular}
  \caption{Analytical pipeline. Each stage tests \textbf{latent routing dimensions} operationalized by prefill statistics---not isolated feature rankings.}
  \label{tab:pipeline}
\end{table}


\paragraph{Statistic definitions.}
\label{sec:method:prefill}
Task difficulty: $c(q)=$ \texttt{piece\_count}.
Model uncertainty and (appendix) model confidence from weak logits at $t{=}T$:
\begin{equation}
H(q,M_i) = -\sum_{v \in V} p_T(v)\log p_T(v), \qquad
m(q,M_i) = p_T^{(1)} - p_T^{(2)},
\end{equation}
full-vocabulary $p_T$, temperature $\tau{=}1$.
Model disagreement and escalation potential:
$\Delta H = H_w - H_s$; $\Delta m_{\mathrm{gain}} = m_s - m_w$.

\paragraph{Branches.}
The same chat prompt $P(u)$ is used for signal extraction and offline oracle labeling; operationalizations are computed offline only.

\begin{table}[t]
  \centering
  \small
  \begin{tabular}{lcc}
    \toprule
    \textbf{Setting} & \textbf{Offline oracle} & \textbf{Prefill probe} \\
    \midrule
    Decoding & Greedy ($\tau{=}0$) & None \\
    Distribution & --- & $\tau{=}1$; full-vocab softmax \\
    Purpose & Labels $y(q,M_i)$ & Logits for operationalizations \\
    \bottomrule
  \end{tabular}
  \caption{Offline oracle labeling vs.\ prefill-based signal extraction.}
  \label{tab:oracle-vs-probe}
\end{table}

\paragraph{Supervision boundary.}
\label{sec:method:supervision}
Signal extraction is unsupervised at inference; dimension-wise characterization uses offline oracle labels as \textbf{evidence}; routing evaluation fits a demonstration policy on CALIB only.

\begin{table}[t]
  \centering
  \small
  \begin{tabular}{@{}p{0.26\linewidth}cp{0.48\linewidth}@{}}
    \toprule
    \textbf{Stage} & \textbf{Labels?} & \textbf{Role} \\
    \midrule
    Signal extraction & No & Unsupervised at inference \\
    Dimension characterization & Oracle offline & \textbf{Evidence} \\
    Routing evaluation & CALIB only & Validate implications \\
    \bottomrule
  \end{tabular}
  \caption{Supervision boundary.}
  \label{tab:supervision}
\end{table}

\subsection{Evaluation design}
\label{sec:method:evaluation}

\textbf{Primary question (Studies I--II):} which latent routing dimension predicts routing opportunity?
\textbf{Secondary question (Study III):} do dimensions provide complementary predictive content?
We report $\rho$, opportunity AUROC/AUPRC, and $\Delta$AUROC when combining operationalized dimensions---not exhaustive leaderboards over every candidate statistic.
Study~III ladder: task difficulty $\rightarrow$ difficulty + model uncertainty $\rightarrow$ joint with model confidence (appendix operationalization).
Routing evaluation (Study~IV) uses operationalizations of task difficulty, model uncertainty, and model confidence.
